\documentclass[10pt,a4paper]{amsart}
\usepackage[margin=19mm]{geometry}
\usepackage{amsmath,amssymb,mathtools}
\usepackage[scr=rsfs,cal=euler]{mathalfa}
\usepackage{xfrac}
\usepackage[unicode,hidelinks,bookmarks=false]{hyperref}
\newcommand{\ie}{i.e.\ }
\newcommand{\cf}{cf.\ }
\newcommand{\resp}{respectively\ }
\newcommand{\Subsection}{\subsection}
\providecommand{\nobreakdash}{\nobreak\hbox{-}}
\setlength{\emergencystretch}{2em}
\multlinegap0pt
\usepackage{url}
\usepackage{stmaryrd}
\usepackage{paralist,mathrsfs}
\usepackage{mathrsfs,booktabs,tabularx}
\usepackage{xifthen,tikz,setspace}
\usetikzlibrary{decorations.pathmorphing,patterns,shapes,calc,decorations}
\usetikzlibrary{decorations.pathreplacing}
\usepackage{mathtools}
\usepackage{todonotes}
\numberwithin{equation}{section}
\renewcommand{\qedsymbol}{\ensuremath{\blacksquare}}
\newcommand{\bin}{\operatorname{Bin}}
\newcommand{\Po}{\operatorname{Po}}
\renewcommand{\restriction}{\mathord{\upharpoonright}}
\renewcommand{\epsilon}{\varepsilon}
\newcommand{\given}{\;\big|\;}
\newcommand{\one}{\mathbf{1}}
\usepackage{color}
\newtheorem{maintheorem}{Theorem}
\newtheorem{thm}{Theorem}
\newtheorem{mainprop}[maintheorem]{Proposition}
\newtheorem{maincoro}[maintheorem]{Corollary}
\newtheorem{conjecture}{Conjecture}
\newtheorem{theorem}{Theorem}[section]
\newtheorem*{theorem*}{Theorem}
\newtheorem{lemma}[theorem]{Lemma}
\newtheorem{lem}[theorem]{Lemma}
\newtheorem{conj}[theorem]{Conjecture}
\newtheorem{claim}[theorem]{Claim}
\newtheorem{proposition}[theorem]{Proposition}
\newtheorem{prop}[theorem]{Proposition}
\newtheorem{observation}[theorem]{Observation}
\newtheorem*{observation*}{Observation}
\newtheorem{fact}[theorem]{Fact}
\newtheorem{corollary}[theorem]{Corollary}
\newtheorem{cor}[theorem]{Corollary}
\newtheorem{assumption}[theorem]{Assumption}
\newtheorem{condition}[theorem]{Condition}
\newtheorem{remark}[theorem]{Remark}
\newtheorem{example}[theorem]{Example}
\newtheorem{definition}[theorem]{Definition}
\newtheorem*{definition*}{Definition}
\newtheorem{problem}[theorem]{Problem}
\newtheorem{question}[theorem]{Question}
\newcommand{\E}{\mathbb E}
\newcommand{\N}{\mathbb N}
\renewcommand{\P}{\mathbb P}
\newcommand{\F}{\mathbb F}
\newcommand{\bbT}{\mathbb T}
\newcommand{\K}{\mathbb K}
\newcommand{\R}{\mathbb R}
\newcommand{\C}{\mathbb C}
\newcommand{\Z}{\mathbb Z}
\newcommand{\cD}{\mathcal D}
\newcommand{\cG}{\mathcal G}
\newcommand{\cE}{\mathcal E}
\newcommand{\cA}{\mathcal A}
\newcommand{\cS}{\mathcal S}
\newcommand{\cF}{\mathcal F}
\newcommand{\cH}{\mathcal H}
\newcommand{\cK}{\mathcal K}
\newcommand{\cT}{\mathcal T}
\newcommand{\cL}{\mathcal L}
\newcommand{\cP}{\mathcal P}
\newcommand{\cV}{\mathcal V}
\newcommand{\cY}{\mathcal Y}
\newcommand{\bC}{\mathbb C}
\newcommand{\fa}{\mathfrak a}
\newcommand{\fb}{\mathfrak b}
\newcommand{\ff}{\mathfrak f}
\newcommand{\fm}{\mathfrak m}
\newcommand{\fs}{\mathfrak s}
\newcommand{\ii}{\mathrm{i}}
\newcommand{\fv}{\mathfrak v}
\newcommand{\sC}{\mathsf C}
\newcommand{\bo}{\mathbf o}
\newcommand{\bu}{\mathbf u}
\newcommand{\bv}{\mathbf r}
\newcommand{\bw}{\mathbf w}
\newcommand{\sL}{\mathscr L}
\newcommand{\p}{{\bf p}}
\newcommand{\q}{{\bf q}}
\newcommand{\tr}{{\mathrm {tr}}}
\newcommand{\Id}{\mathrm{Id}}
\newcommand{\Bin}{\mathrm{Bin}}
\newcommand{\aG}{\fa}
\newcommand{\aGhat}{\widehat\fa}
\newcommand{\bG}{\fb}
\newcommand{\leqP}{\mathbin{\leq_{\textsc p}}}
\newcommand{\nleqP}{\mathbin{\nleq_{\textsc p}}}
\DeclareMathOperator{\var}{Var}
\DeclareMathOperator{\dist}{dist}
\DeclareMathOperator{\diam}{diam}
\DeclareMathOperator{\cov}{Cov}
\DeclareMathOperator{\rank}{rank}
\begin{document}
\title[On the d-rigidity phase transition in random graphs]{On the $d$-rigidity phase transition in random graphs}
\author{Yuval Peled}
\date{}
\maketitle
\noindent\emph{Source and licence.} On the $d$-rigidity phase transition in random graphs. Version arXiv:2605.25711v3 (CC BY 4.0). This material is distributed under the Creative Commons Attribution 4.0 International licence, http://creativecommons.org/licenses/by/4.0/, as recorded in the arXiv OAI licence metadata for this exact version and shown on the abstract page https://arxiv.org/abs/2605.25711v3. Full rights notice: licenses/arxiv-2605.25711-CC-BY-4.0.txt.\par\smallskip
\noindent\textit{Editorial scope.} Counted unit: the original \texttt{theorem} environment \texttt{thm:recursion} at source lines 907--926 together with its complete proof at lines 928--1060; the delivered document keeps source lines 454--1281 so that every referenced label and every citation resolves inside it. Native editable source is the paper's own concatenated TeX; the AMS wrapper is editorial formatting only. No publisher class, upstream script or unrelated artwork is redistributed.
\par\smallskip\noindent\textit{Reference note.} This article ships no compiled bibliography (biblatex); the entries delivered here are composed from the author's own BibTeX (.bib) fields, with no field invented and no entry dropped.
\section{Preliminaries}\label{sec:preliminaries}
\subsection{Infinitesimal rigidity theory}
Throughout this paper, we fix a dimension $d\ge 2$. Let $G$ be an $n$-vertex graph and $\p:V(G)\to\R^d$ a generic embedding. We view every $\p(v)$ as a column vector in $\R^d$, 
and denote $\p_{uv}=\p(u)-\p(v)$ for two distinct vertices $u$ and $v$. 
The rigidity matrix $R = R(G,\p)$ of the $d$-framework $(G,\p)$ is the $|E| \times d|V|$ matrix whose rows are indexed by edges $uv \in E$ and whose columns are grouped into $|V|$ blocks of size $d$, one block for each vertex. 
For an edge $uv \in E$, the row corresponding to $uv$ has $\p_{uv}^*$ in the block corresponding to $u$, $\p_{vu}^*$ in the block corresponding to $v$, and zero elsewhere.

Throughout this paper it will be more convenient to work with the {\em rigidity Laplacian} $L=L(G,\p)=R^TR$ (defined more explicitly below), as it is a self-adjoint operator.
We denote by $Z=Z(G,\p)$ the kernel of $R(G,\p)$; equivalently, $Z$ is the kernel of the rigidity Laplacian $L(G,\p)$. 
A \(dn\)-dimensional vector \(q\) in \(Z\) can be viewed as assigning an
infinitesimal velocity vector \(q(v)\in\R^d\) to each vertex, describing
an infinitesimal motion of the vertices at \(\p\) under which all edge
lengths are preserved to first order. If \(\p\) is generic then the
isometries of \(\R^d\) induce a \(\binom{d+1}{2}\)-dimensional subspace of
\(Z\), which we refer to as the {\em trivial motions}. More concretely, the
subspace of trivial motions is parametrized by a \(d\times d\)
skew-symmetric matrix \(A\) and a vector \(t\in\R^d\), yielding a vector
\(q\in Z\) defined by
\[
q(v)=A\p(v)+t,\qquad v\in V.
\]
$G$ is $d$-rigid iff $Z$ consists only of trivial motions. 

Another matrix that will play a role in this work is the orthogonal
projection matrix \(P:\R^{dn}\to Z\). Intuitively, suppose that each vertex
has unit mass and that the edges are massless bars. If \(f\in \R^{dn}\) describes a vector of instantaneous external forces,
where \(f(v)\in\R^d\) is the force applied to the vertex \(v\), then
\(Pf\) is the component of this force in the space of admissible
infinitesimal motions; equivalently, for unit masses, it is the resulting
instantaneous acceleration after imposing the bar constraints.
One may consider this locally. Let \(P_{v,v}\) denote the
\(d\times d\) diagonal block of \(P\) corresponding to the vertex \(v\).
If a force \(f(v)\) is applied only to \(v\), then the resulting
instantaneous acceleration of \(v\) is
$P_{v,v} f(v).$ Thus, the trace of $P_{v,v}$ quantifies the degrees of freedom that $v$ has in the framework, and we will refer to it as the local flexibility of $v$; see Section \ref{sec:local_flex} below for a formal definition and further discussion.

\subsection{Benjamini--Schramm convergence}
A main idea in this work is to estimate the rigidity rank of a finite graph by its local structure. 
For this, we work with the local weak topology of rooted graphs introduced in the influential paper 
of Benjamini and Schramm \cite{BenjaminiSchramm2001}. 
This point of view is closely related to the objective method of Aldous and Steele
\cite{AldousSteeleObjectiveMethod}, and has been used to study a variety of
parameters of sparse graphs which are governed by local limits. Of particular
relevance to us are results concerning spectral measures and kernels of
operators on locally convergent graphs; see, for example,
\cite{AbertThomViragSpectralMeasure,BordenaveLelarge2010,BordenaveLelargeSalezDiluted}.

A rooted graph $(G,o)$ is a graph with a root vertex $o\in V(G)$. Let $\cG_\bullet$ denote the set of isomorphism classes of connected locally finite rooted graphs $(G,o)$, 
where isomorphisms are required to preserve the root. 
For $(G,o)\in\cG_\bullet$ and an integer $r\ge 0$, we write $B_G(o,r)$ for the rooted induced subgraph spanned by all vertices 
at graph distance at most $r$ from $o$; thus $B_G(o,r)\in\cG_\bullet$ as well. We metrize $\cG_\bullet$ by defining 
the distance between $(G_1,o_1)$ and $(G_2,o_2)$ to be $2^{-r}$, where $r\ge 0$ is the largest radius for which the 
rooted balls $B_{G_1}(o_1,r)$ and $B_{G_2}(o_2,r)$ are root-preserving isomorphic. The induced metric is the 
{\em local topology} on $\cG_\bullet$; with this topology, $\cG_\bullet$ is a Polish space, so the standard notions of tightness and convergence in distribution apply.
We say that a (possibly random) rooted graph $(G,o)$ is the local weak (Benjamini--Schramm) limit of a sequence of (possibly random) finite graphs $G_n$ if $(G_n,o_n)$ 
converges in distribution to $(G,o)$ in $\cG_\bullet$, where $o_n$ is chosen uniformly from $V(G_n)$.

We mention two classical examples of local weak convergence from graph theory that appear in this paper. First, $G(n,c/n)$, where $c>0$ is fixed, converges locally to a 
Galton--Watson branching process with $\operatorname{Poi}(c)$ progeny distribution (defined precisely in the next subsection) as $n\to\infty$. 
Second, the uniform random $k$-regular $n$-vertex graph, where $k\ge 3$ is fixed, converges locally to the infinite $k$-regular tree $\bbT_k$ as $n\to\infty$. 
Both these examples can be viewed as special cases of local weak convergence of uniform random graphs with a given degree distribution 
to a corresponding unimodular Galton--Watson process.

\subsection{Galton--Watson processes, size-biasing and thinning}
The graphs we study here converge locally to Galton--Watson processes. 
A Galton--Watson branching process is a random rooted tree generated recursively by letting each vertex produce an independent number of children.  
Let $X,Y$ be probability distributions on $\N_{\ge 0}$. Denote by $\mathrm{GW}(X,Y)$ the distribution on random, locally finite rooted trees $T$ generated 
by the Galton--Watson branching process in which the root $o$ has offspring distribution $X$, while every other vertex has offspring distribution $Y$. 

Galton--Watson processes that arise as local limits of finite graphs are unimodular, which is equivalent to the condition that $Y$ is the size-biasing of $X$:
\begin{equation}\label{eq:size_bias}
Y(m-1)=\frac{m\,X(m)}{\E[X]}\,,\quad\forall m\ge 1\,.
\end{equation}
Here we assume $\E[X]<\infty$.

Two main examples that we consider in this paper are:
\begin{enumerate}
    \item \(X=Y=\operatorname{Poi}(c)\), for some fixed \(c>0\), which is the
    local weak limit of \(G(n,c/n)\);
    \item the \(k\)-regular tree \(\bbT_k\), i.e.,
    \(X(k)=Y(k-1)=1\), which is the local weak limit of the random
    \(k\)-regular graph.
\end{enumerate}


An easy and fundamental fact about size-biasing is that for every $f:\N_{\ge 1}\to\R$ there holds,
\begin{equation}\label{eq:sizeBiasF}
\E[Xf(X-1)] = \E[X]\E[f(Y)]\,.
\end{equation}

In addition, we will use the operation of thinning, and its relation to size-biasing. Given $0\le p \le 1$, we define a coupled pair $(X_p,X_{1-p})$ by 
first sampling $X$, and then, conditional on $X$, sampling $X_p\sim \Bin(X,p)$ and setting $X_{1-p}:=X-X_p$. We can also view this pair as coupled with $X$.
The following claims are useful in our argument.
\begin{claim}\label{clm:sizeBias_new_claim}
Suppose that $Y$ is the size-biasing of $X$ and $0\le p\le 1$. Then,
\begin{enumerate}
\item $Y_p$ is the size-biasing of $X_p$.
\item For every positive integer $d>0$ there holds $$(1-p)\E[X]\P(Y_p<d)=\E[\mathbf 1_{X_p<d}\cdot X_{1-p}]\,.$$
\item Suppose that $t\sim Y_{1-p}$ and that conditioned on $t$, $h\sim X_p | \{X_{1-p}=t+1\}$. Then, the pair $(h,t)$ is $(Y_p,Y_{1-p})$-distributed.
\end{enumerate}
\end{claim}
\begin{proof}
First, for every $t \ge 0$ there holds 
\begin{align}
\label{eq:yp_sizebias_of_xp}
(t+1)\,X_p(t+1)
&= \E\!\left[\mathbf 1_{X_p=t+1}\cdot X_p\right] \notag\\
&= \E\!\left[\E\!\left[\mathbf 1_{X_p=t+1}\cdot X_p \mid X\right]\right] \notag\\
&= \E\!\left[X\cdot p \cdot \P(\Bin(X-1,p)=t )\right] \notag\\
&= p\,\E[X]\cdot \E[\P(\Bin(Y,p)=t)] \notag\\
&= \E[X_p]\cdot Y_p(t)\,, \notag
\end{align}
as claimed. The third equality follows from linearity of expectation, since given $X$, $\mathbf 1_{X_p=t+1}\cdot X_p$ is a sum of $X$ Bernoullis with probability $p\cdot\P(\Bin(X-1,p)=t).$ In addition, the fourth equality is derived by \eqref{eq:sizeBiasF} using $f(k)=\P(\Bin(k,p)=t).$


Similarly, for the second item, by linearity of expectation,
\[
\E[\mathbf 1_{X_p<d}\cdot X_{1-p}\mid X] = X\cdot (1-p)\cdot\P(\Bin(X-1,p)<d)\,,
\]
and by taking expectation over $X$ we obtain
\begin{align*}
\E\!\left[\mathbf 1_{X_p<d}\,X_{1-p}\right]
&=(1-p)\,\E\!\left[X\cdot\P\!\left(\Bin(X-1,p)<d \right)\right] \\
&=(1-p)\,\E[X]\cdot\E[\P\!\left(\Bin(Y,p)<d\right)]
\;=\;(1-p)\,\E[X]\,\P\!\left(Y_p<d\right).
\end{align*}
The transition to the second line is derived by \eqref{eq:sizeBiasF} using $f(k)=\P(\Bin(k,p)<d)$.

For the third item, note that for every $t,h$ it holds that
\begin{align*}
\P(X_p=h\mid X_{1-p}=t+1) &=~ \frac{X(h+t+1)\P(\Bin(h+t+1,p)=h)}{X_{1-p}(t+1)}\\[0.6em]
&=~ \frac{(1-p)(h+t+1)X(h+t+1)\P(\Bin(h+t,p)=h)}{(t+1)X_{1-p}(t+1)} \\[0.6em]
&=~ \frac{(1-p)\E[X]Y(h+t)\P(\Bin(h+t,p)=h)}{(t+1)X_{1-p}(t+1)} \\[0.6em]
&=~ \frac{Y(h+t)\P(\Bin(h+t,p)=h)}{Y_{1-p}(t)} \\
&=~ \P(Y_p=h\mid Y_{1-p}=t) \,,
\end{align*}
and the claim follows since we assume $t\sim Y_{1-p}$. Here, the second equality follows from a standard identity of Binomial distributions, the next equality holds since $Y$ is the size-biasing of $X$, and the fourth equality is derived by the first item.
\end{proof}


\section{Infinite frameworks}\label{sec:inf_fram}
An important idea in this work is to study rigidity of infinite graphs, and
infinite trees in particular. Rigidity of infinite frameworks was studied by
Owen and Power \cite{OwenPowerInfiniteFrameworksOperatorTheory} and by
Kitson and Power \cite{KitsonPowerRigidityInfiniteGraphs}. The approach we
take here is closest to the notion of square-summable
infinitesimal rigidity from \cite{OwenPowerInfiniteFrameworksOperatorTheory}.
However, we opt to give a self-contained treatment, partly since our focus is on the spectral-theoretic aspects of this notion.

The infinite graphs we consider are locally finite but may have unbounded degree. Consequently, their rigidity Laplacians may be unbounded operators. Accordingly, we begin by recalling some background material from the theory of unbounded operators~\cite{ReedSimonI}.
Let $\cH$ be a complex Hilbert space. The main issue with unbounded operators is that they cannot be defined on the entire space $\cH$. Hence, an unbounded operator on $\cH$ consists of a domain $D$ which is a linear subspace of $\cH$, and a linear operator $L:D\to \cH$. As we aim to apply spectral methods, we need to work with a self-adjoint operator. However, if $L$ is not closed, that is, if the graph $\{(f,Lf)\mid f\in D\}$ is not closed in $\cH\oplus\cH$, then $L$ cannot be self-adjoint, since the adjoint $L^*$ is always closed. 
Fortunately, in our setting, $D$ is a {\em dense} subspace and $L$ is symmetric. Therefore, $L$ is closable, and we denote its closure by $\bar L$. The domain $D$ is a {\em core} for $\bar L$.
We say that $L$ is essentially self-adjoint if $\bar L$ is self-adjoint.
In this case, the spectral theorem applies to $\bar L$. In particular, for any bounded Borel function $g:\R\to\C$, the operator $g(\bar L)$ is well defined via the functional calculus. Moreover, for every vector $f\in\cH$, there exists a finite Borel measure $\mu_f$ on $\R$, called the spectral measure of $\bar L$ with respect to $f$, such that
\[
\langle f, g(\bar L) f\rangle = \int_\R g(\lambda)\, d\mu_f(\lambda)
\]
for all bounded Borel functions $g$. In particular, if $g=\mathbf 1_{\{0\}}$ then $g(\bar L)$ is the orthogonal projection $P:\cH\to\ker\bar L$ and $\langle f,Pf\rangle=\mu_f(\{0\})$.

All the graphs $G=(V,E)$ we consider
are locally finite on a countable vertex set. Let $\p:V\to\R^d$. Since we use spectral theory, we prefer to work over the complex numbers and consider the Hilbert space
\[
\cH:=\ell_2(V;\C^d)\,.
\]
This makes no difference to us since rigidity theories over $\R$ and $\C$ are identical. Denote by $D\subset\cH$ the subspace of finitely supported vectors. The rigidity Laplacian $L=L(G,\p)$ is given by a (possibly infinite) block matrix
\[
L=(L_{uv})_{u,v\in V}, \qquad 
L_{uv}\in\C^{d\times d},
\]
acting on $f\in D$ by
\[
(L f)(u)=\sum_{v\sim u}L_{uv} f(v),
\]
where the sum is finite for each $u$. The block entries are given by
\begin{equation}\label{eq:Luv}
L_{uv}=
\begin{cases}
\displaystyle \sum_{w\sim u} \p_{uw}\p_{uw}^{*}, & u=v,\\[1ex]
\displaystyle -\,\p_{uv}\p_{uv}^{*}, & \{u,v\}\in E,\\[1ex]
0, & \text{otherwise}\,.
\end{cases}
\end{equation}


The operator $L$ is symmetric and defined on the dense domain $D$, since $L_{uv}=L_{vu}^*$ for all $u,v\in V$. Therefore, $L$ is closable, and we denote its closure by $\bar L$.

\begin{definition}
Let $G$ be a locally finite graph on a countable vertex set and $\p:V(G)\to\R^d$.
The $d$-framework $(G,\p)$ is called \emph{self-adjoint} if $\bar L(G,\p)$ is a self-adjoint operator.
\end{definition}
It is clear that if $G$ is finite, or, more generally, if $G$ has bounded degree and $\p_{uv},~uv\in E$ are bounded, then $(G,\p)$ is self-adjoint. Indeed, in such a case the Rigidity Laplacian is a bounded operator.

Although self-adjointness of a framework is defined via the complex Hilbert
space \(\cH\), if a framework \((G,\p)\) is self-adjoint, we note that the
orthogonal projection
\[
  P=\mathbf 1_{\{0\}}(\bar L):\cH\to \ker(\bar L)
\]
onto the kernel of its rigidity Laplacian can be viewed as a real operator. Indeed, the
rigidity Laplacian has real entries, hence the projection \(P\) preserves the real subspace of $\cH$.
Thus, when applied to real vectors, \(P\) is the orthogonal projection onto
the real kernel of \(\bar L\). 


\subsection{$\ell_2$-infinitesimal motions}
The following lemma extends the standard description for the kernel of the rigidity Laplacian from finite graphs to self-adjoint frameworks. In words, the kernel consists of all infinitesimal distance preserving motions that have a finite $\ell_2$-norm.
\begin{lemma}\label{lem:kernel}
Let $(G,\p)$ be a self-adjoint $d$-framework. Then,
\[
\ker \bar L(G,\p) = \{q\in\cH\mid \langle q(u)-q(v),\p_{uv}\rangle=0,~~\forall uv \in E(G)\}\,.
\]
\end{lemma}
\begin{proof}
The dense subspace $D$ of finitely-supported vectors is a {\em core} for the self-adjoint operator $\bar L=\bar L(G,\p)$. That is, for every $q$ in the domain of $\bar L$ there are vectors $f_n \in D$ such that $f_n\to q$ and $Lf_n\to \bar L q$. 

Suppose that $q\in \ker \bar L$, and consider vectors $f_n \in D$ such that $f_n\to q$ and $Lf_n\to 0$. We have that
\[
\sum_{uv\in E}|\langle f_n(u)-f_n(v),\p_{uv}\rangle|^2 = \langle f_n,Lf_n\rangle \to \langle q,\bar Lq\rangle =0\,.
\]
Hence, for every edge $uv \in E$, 
\[
\langle q(u)-q(v),\p_{uv}\rangle=\lim_{n\to \infty}\langle f_n(u)-f_n(v),\p_{uv}\rangle = 0\,.
\]

On the other hand, suppose that $q\in\cH$ satisfies $\langle q(u)-q(v),\p_{uv}\rangle=0$ for every edge $uv\in E$. Let $f\in D$. Then,
\begin{align*}
\langle q,Lf\rangle &= \sum_{u \in V}\left\langle q(u),\sum_{v\sim u}\p_{uv}\p_{uv}^*(f(u)-f(v)) \right\rangle \\
&= \sum_{uv\in E}\left\langle q(u)-q(v),\p_{uv}\p_{uv}^*(f(u)-f(v)) \right\rangle=0\,,
\end{align*}
where the last equality follows from our assumption on $q$. Note that this summation is finite since $f\in D$. Therefore, by the definition of the adjoint, $\bar L^*q=0$, and the proof follows by the assumption that $\bar L$ is self-adjoint.
\end{proof}

The following definition is natural, but will not be of much use in this work.
\begin{definition}
    A self-adjoint $d$-framework $(G,\p)$ is called $\ell_2$ \emph{strictly $d$-rigid} if \[\ker \bar L(G,\p)=\{0\}\,.\]
\end{definition}

No finite graph is strictly-rigid since \(\ker L\) contains the trivial motions.
For infinite graphs, $\ell_2$ strict rigidity is more subtle since the trivial motions need not belong to
\(\ell_2\). Consequently, \(\ker \bar L\) may be trivial even when
\(\p\) is generic, as the following example illustrates. The example
also shows that, on the other hand, there are generic
embeddings for which \(\ker L\) is non-trivial even for infinite graphs.

\begin{example}\label{ex:strictly_rigid}
Let $G$ be a bounded-degree graph on $V=\N$ such that the induced subgraph $G_n$ on
$[n]=\{1,\dots,n\}$ is $d$-rigid for every $n\in\N$. Let $\p:V\to[0,1]^d$ be a generic embedding which is \emph{well-spread}, in the sense that for some $\varepsilon>0$, and for every affine subspace $F\subset\R^d$, there exist infinitely many $v\in V$ such that
\[
\dist(\p(v),F)\ge \varepsilon.
\]
For example, draw $\p(v)$ uniformly and independently from $[0,1]^d$ for each $v\in V$.
By boundedness of $G$ and $\p$, the framework $(G,\p)$ is self-adjoint. Let $q\in\ker(\bar L)$. Restricting $q$ to the induced subgraph $G_n$, the
assumption that $G_n$ is $d$-rigid implies that $q|_{[n]}$ is induced by a trivial motion. Since this holds for every $n$, it follows from the genericity of $\p$ that $q$ is induced by a single global trivial motion. That is, there exist a skew-symmetric matrix $A$ and a vector $t$ such that
\[
q(v)=A\p(v)+t \qquad \text{for all } v\in V.
\]
The well-spreadness assumption on $\p$ implies that
\begin{equation}\label{eq:series}
\sum_{v\in V}\|A\p(v)+t\|^2 = \infty
\end{equation}
for all $(A,t)\neq(0,0)$. Since $q\in\cH=\ell_2(V;\C^d)$, this forces $A=0$ and $t=0$, and hence $q=0$. Therefore,
$
\ker(\bar L)=\{0\}.
$

Note that a different choice of a (possibly) generic embedding $\p$ may lead to non-zero elements of $\cH$ arising from trivial motions: indeed, if there exist a skew-symmetric matrix $A$ and a vector $t\in\R^d$, not both zero, such that \eqref{eq:series} converges, then the corresponding trivial motion belongs to $\cH$ and hence lies in $\ker(\bar L)$.
\end{example}

In the example above, every vertex is completely constrained: no vertex
can move in any direction. However, this phenomenon can be localized to
some directions but not others. For instance, if we add a new vertex \(v\)
and connect it to one vertex \(u\) of the original graph, then the motion
of \(v\) is constrained to lie in the subspace orthogonal to \(\p_{uv}\).
As before, such behavior cannot occur in finite frameworks, because
trivial motions are always present. The following lemma characterizes the
forbidden directions for a vertex.

\begin{lemma}\label{lem:forbidden_directions}
    Let $(G,\p)$ be a self-adjoint $d$-framework, and $v\in V(G)$ a vertex. Let $P:\cH\to\ker\bar L(G,\p)$ 
    be the orthogonal projection, and $P_{v,v}$ the $d\times d$ block corresponding to $v$. Then, for every $x\in \C^d$,
    \[
    x\in\ker P_{v,v} \iff \langle x,q(v)\rangle=0,~~\forall q\in\ker \bar L(G,\p)\,.
    \]
\end{lemma}
\begin{proof}
    Let $f\in D$ be defined by $f(v)=x$ and $f(u)=0$ for every $u\ne v$. Then,
    \[
    \|Pf\|^2 = \langle f,Pf\rangle = \langle x,(Pf)(v)\rangle = \langle x,P_{v,v}x\rangle\,.
    \]
    Therefore, since $P_{v,v,}$ is PSD, $P_{v,v}x=0$ if and only if $Pf=0$. Since $P$ is an orthogonal projection, this is equivalent to the condition that for every $q\in\ker \bar L(G,\p),$
    \[
    \langle x,q(v)\rangle = \langle f, q\rangle = 0\,,
    \]
    as claimed.
\end{proof}

\subsection{Random frameworks}\label{sec:admissible}
Due to the illustrated subtleties in generic rigidity of infinite frameworks, 
and because it is better suited to our methods, we shall often restrict the class of embeddings considered in this work as follows. 
Given a possibly random graph \(G\) and a probability measure \(\nu\) on
\(\R^d\), we write \((G,\nu)\) for the random \(d\)-framework \((G,\p)\)
in which the vectors \(\p(v)\), \(v\in V(G)\), are sampled independently
from \(\nu\).

We shall mostly consider probability measures \(\nu\) on \(\R^d\) that we
call \emph{admissible}: namely, \(\nu\) is supported on a bounded nonempty
open set \(U\subset\R^d\), and is absolutely continuous with respect to
Lebesgue measure on \(U\). For example, one may take \(\nu\) to be the
uniform distribution on \([0,1]^d\).


\section{Infinitesimal local flexibility}\label{sec:local_flex}
We now introduce the notion of local flexibility, which plays a central
role in this work. Although the notion can be defined for finite
frameworks, we formulate it in a more general setting.

Let $(G,\p)$ be a self-adjoint $d$-framework. 
Recall that 
    \[
    P=\mathbf 1_{\{0\}}(\bar L):\cH\to \ker(\bar L)
    \]
    and $P_{v,v}$ is the real $d\times d$ block of $P$ corresponding to a vertex $v$. 
\begin{definition}
Let $(G,\p)$ be a self-adjoint $d$-framework and $v\in V(G)$. The {\em local flexibility} of $v$ in $(G,\p)$ is defined by
    \[
    \phi_{G,\p}(v)=\tr(P_{v,v})\,.
    \]
\end{definition}
Observe that an alternative description of the local flexibility is given by
\[
\phi_{G,\p}(v) = \sum_{j=1}^d \mu_{e_{v,j}}(\{0\})\,,
\]
where $e_{v,j}$ is the unit vector of the $j$-th coordinate in the block corresponding to $v$, and $\mu_{e_{v,j}}$ is the spectral measure of $\bar L$ with respect to $e_{v,j}$.

When the context is clear, we sometimes omit the subscript $G,\p$ in $\phi$.

Intuitively, the local flexibility quantifies the degrees of freedom that $v$ has in the framework. 
In particular, it is easy to see that $d-\deg_G(v)\le \phi (v)\le d$. In a finite graph, $\phi(v)>0$ due to the trivial motions. 
But, Example \ref{ex:strictly_rigid} shows that the local flexibility can vanish in infinite graphs.


Crucially, by standard linear algebra, if $G$ is finite then the dimension of the kernel of the rigidity matrix $R$ is given by
\begin{equation}\label{eq:dimZ_local}
\dim Z(G,\p) = \sum_{v\in V(G)} \phi_{G,\p}(v)\,,
\end{equation}
since both sides are equal to the trace of $P$. Thus we  view the local flexibility of $v$ as the local ``contribution'' of $v$ to the dimension $Z$.

\subsection{Upper semi-continuity under Benjamini--Schramm convergence}
We now show that the local flexibility $\phi_{G,\p}(v)$ is a {\em local parameter}. That is, it is upper semi-continuous 
under Benjamini--Schramm convergence. This is a crucial property for our applications, as it allows us to deduce results about the local flexibility of finite graphs from results about the local flexibility of their infinite local limit.

This is a direct adaptation to the rigidity Laplacian of previous results on the weak convergence of spectral measures of graph operators under local weak convergence (e.g., \cite{BordenaveLelargeSalezDiluted,AbertThomViragSpectralMeasure}).

\begin{lemma}\label{lem:upper_semi}
    Suppose that $(G,o)$ is the local weak limit of a graph sequence $G_n$. Let $\nu$ be a probability measure on $\R^d$ for which $(G,\nu)$ and $(G_n,\nu)$ are all almost surely self-adjoint. Then,
    \[
        \limsup_{n\to\infty}\E[\phi_{G_n,\nu}(o_n)]\le \E[\phi_{G,\nu}(o)]\,,
    \]
    where $o_n\in V(G_n)$ is a uniform random vertex.
\end{lemma}
\begin{proof}
     By the Skorokhod representation theorem, we may assume that $(G_n,o_n)\to (G,o)$ almost surely in $\cG_\bullet$. By choosing appropriate embeddings of all graphs into a common vertex set $V$, we may further assume that $o_n=o\in V$ for all $n$, and that for every $r\ge 2$ there exists $n_r$ such that, almost surely,
    \begin{equation}\label{eq:equal_balls}
    B_{G_n}(o,r)=B_G(o,r)\qquad\text{for all } n\ge n_r.
    \end{equation}
    Fix $j\in[d]$ and let $e_j = e_{o,j}\in\ell_2(V;\C^d)$ be the unit vector corresponding to the $j$-th coordinate at the root $o$.

    Let $\p:V\to \R^d$ be such that $\p(v)$, $v\in V(G)$, are independent samples of $\nu$, and denote $\p_n=\p|_{V(G_n)}$.
    By assumption, the frameworks $(G_n,\p_n)$ and $(G,\p)$ are almost surely self-adjoint. Hence the associated rigidity Laplacians $\bar L_n = \bar L(G_n,\p_n)$ and $\bar L = \bar L(G,\p)$ admit spectral measures $\mu^{(n)}_{e_{j}}$ and $\mu_{e_{j}}$ with respect to $e_{j}$.

From \eqref{eq:equal_balls} and the assumption $\p_n=\p|_{V(G_n)}$, we deduce that for every \(f\in D\), \(L_n f=Lf\) for all sufficiently large \(n\),
since $L_n,L$ are equal in the blocks of the support of $f$. 
Due to the fact that \(D\) is a core for all \(\bar L_n\) and \(\bar L\), \cite[Theorem VIII.25(a)]{ReedSimonI} implies that \(\bar L_n\to \bar L\) in the strong resolvent sense. Therefore, we obtain almost sure weak convergence of the spectral measures $\mu^{(n)}_{e_{j}}\Rightarrow \mu_{e_{j}}$. By the Portmanteau theorem,
    \[
        \limsup_{n\to\infty}\mu^{(n)}_{e_{j}}(\{0\})\le \mu_{e_{j}}(\{0\})\,.
    \]
    Summing over $j=1,\dots,d$ and recalling that $\phi_{G,\p}(o)=\sum_{j=1}^d \mu_{e_{o,j}}(\{0\})$, we obtain, almost surely,
    \[
        \limsup_{n\to\infty}\phi_{G_n,\p_n}(o_n)\le \phi_{G,\p}(o)\,.
    \]
    The proof then follows by Fatou's lemma, using that $0\le \phi_{G_n,\p_n}(o_n)\le d$.
\end{proof}

\section{Local flexibility in trees}\label{sec:trees}
\subsection{Self-adjointness of trees}
From a classical rigidity-theory perspective, finite trees are fairly easy to analyze and are arguably not particularly interesting. However, as we see here, from the perspective of local flexibility, trees---and infinite trees in particular---are nontrivial and very interesting to study. 

We start this investigation by formulating a useful sufficient condition for self-adjointness of frameworks whose underlying graph 
is an infinite tree $T$. We say that a finite subset $S\subset V(T)$ that induces a connected subtree is $C$-bounded, for some $C>0$, if no vertex in $S$ has more than $C$ neighbors outside of $S$.
\begin{lemma}\label{lem:tree_self_adjoint}
    Let $(T,o)$ be a locally finite rooted tree and $\p:V(T)\to\R^d$ such that $\{\p_{uv}\mid uv\in E\}$ is bounded. If there exists $C>0$ such that $V(T)$ is covered by $C$-bounded subsets containing $o$, then $(T,\p)$ is self-adjoint.
\end{lemma} 
\begin{proof}
    By \cite[Theorem VIII.3]{ReedSimonI}, to show that $\bar L=\bar L(T,\p)$ is self-adjoint it suffices to prove that there is no non-zero vector $\psi$ in the domain of $\bar L^*$ for which $L^*\psi=\pm \ii\psi$. We consider the case $\bar L^*\psi=-\ii\psi$, and note that the proof of the other case is nearly identical. For a vertex $u$, let $\psi_u\in D$ be defined by $\psi_u(u)=\psi(u)$ and $\psi_u(v)=0$ for every $v\ne u$. Therefore,
    \[
        \bar L\psi_u(v) = L\psi_u(v) =
        \begin{cases}
        \displaystyle \sum_{w\sim u} \langle \psi(u),\p_{uw}\rangle\p_{uw}, & u=v,\\[1ex]
        \displaystyle -\,\langle \psi(u),\p_{uv}\rangle\p_{uv}, & \{u,v\}\in E,\\[1ex]
        0, & \text{otherwise}\,.
        \end{cases}
    \]
    Hence,
    \[
        \langle \bar L\psi_u,\psi\rangle =\sum_{v\sim u}\langle\psi(u),\p_{uv}\rangle\langle\p_{uv},\psi(u)-\psi(v)\rangle\,. 
    \]
    On the other hand, by our assumption on $\psi$,
    \[
        \langle \bar L\psi_u,\psi\rangle = \langle \psi_u,\bar L^*\psi \rangle = +\ii\|\psi(u)\|^2\,.
    \]    
    Therefore, for any finite subset $S\subset V(T)$ there holds
    \[
        \ii\sum_{u\in S}\|\psi(u)\|^2 =
        \sum_{uv\in E(S)}|\langle\psi(u)-\psi(v),\p_{uv}\rangle|^2
        +\sum_{uv\in E(S,S^c)}\langle\psi(u),\p_{uv}\rangle\langle\p_{uv},\psi(u)-\psi(v)\rangle\,. 
        \]
        Here $E(S)$ denotes the edges induced by subset $S$, and $uv\in E(S,S^c)$ denotes that $u\in S,v\notin S$ and $u\sim v$. By considering the imaginary part of this equality and applying the triangle and Cauchy--Schwarz (CS) inequalities we obtain 
        \begin{equation}\label{eq:sa_norm}
        \sum_{u\in S}\|\psi(u)\|^2 \leq
        \sum_{uv\in E(S,S^c)} \|\psi(u)\|\cdot \|\p_{uv}\|^2\cdot \|\psi(v)\|\,.
        \end{equation}
        Suppose that $S$ is $C$-bounded, and that $\|\p_{uv}\|^2\le C$ for all $uv\in E$. In such a case, we find by applying CS twice that
        \begin{align*}
            \sum_{u\in S}\|\psi(u)\|^2 
            \leq ~&C\cdot \sum_{u\in S} \|\psi(u)\|\cdot C^{1/2}\cdot\left(\sum_{v\sim u,v\notin S}\|\psi(v)\|^2\right)^{1/2}\\
            \leq~
            &C^{1.5}\cdot \left(\sum_{u\in S} \|\psi(u)\|^2\right)^{1/2}\cdot\left(\sum_{u\in S}\sum_{v\sim u,v\notin S}\|\psi(v)\|^2\right)^{1/2} \\
            \leq~
            & C^{1.5}\cdot \left(\sum_{u\in S} \|\psi(u)\|^2\right)^{1/2}\cdot\left(\sum_{v\notin S}\|\psi(v)\|^2\right)^{1/2}\,.
        \end{align*}
    The last inequality follows from the assumption that $S$ induces a connected subtree, hence every $v\notin S$ has at most one neighbor in $S$. A straightforward manipulation yields
        \[
        \sum_{u\in S}\|\psi(u)\|^2 \leq \frac{C^3}{C^3+1}\|\psi\|^2\,.   
        \]
    On the other hand, the family of $C$-bounded subsets that contain the root is closed under taking finite unions, whence our assumption that $V(T)$ is covered by such sets implies that the supremum of $\sum_{u\in S}\|\psi(u)\|^2$ over all $C$-bounded subsets $S$ is $\|\psi\|^2$, yielding a contradiction.
\end{proof}

\subsection{Forward flexibility matrices and their recursive formula}\label{sec:recursive}

For a rooted tree $(T,o)$ and a vertex $v\in V(T)$ we denote by $T^v$ the subtree of $T$ induced by $v$ and its descendants --- all the vertices $u$ for which the shortest path in $T$ from $o$ to $u$ passes through $v$ --- 
and view it as rooted at $v$. 
In addition, for an embedding $\p:V\to \R^d$, we denote by $\p^v=\p|_{T^v}$. 

It is easy to see that if $(T,\p)$ is self-adjoint then so is $(T^v,\p^v)$. In such a case, we consider the orthogonal projection
\[
P^v=\mathbf 1_{\{0\}}(\bar L(T^v,\p^v)):\ell_2(V(T^v);\C^d)\to\ker(\bar L(T^v,\p^v))\,,
\]
and call its restriction $Q_v=(P^v)_{vv}$ to the $d\times d$ block corresponding to the vertex $v$ the {\em forward flexibility matrix} of $v$ in $T$.
Note that $Q_v$ has real entries, and recall that its trace is the local flexibility of $v$ in $(T^v,\p^v)$. 
In addition, by Lemma \ref{lem:forbidden_directions}, the kernel of $Q_v$ consists of the forbidden directions for $v$ in this framework $(T^v,\p^v)$.

The next theorem finds a recursive formula for computing $Q_o$ using the vectors $\p_{ov}$ and the matrices $Q_v$, over all children $v$ of the root $o$.


\begin{theorem}\label{thm:recursion}
Let $(T,\p)$ be a self-adjoint $d$-framework whose underlying graph $T$ is a tree rooted at $o$. Denote by
$v_1,\dots,v_m$ the children of the root $o$. Let
\[
H=\{\,i\in[m]\mid Q_{v_i}\p_{ov_i}=0\,\},
\]
$W\le \C^d$ be the orthogonal complement of
$\mathrm{span}_{\C}\{\p_{ov_i}\mid i\in H\}$, and $P^W:\C^d\to W$ be the orthogonal projection. Then,
\begin{enumerate}
\item $Q_o\p_{ov_i}=0$ for every $i\in H$.
\item
\[
Q_o|_W
=
\left(
\Id_W+\sum_{i\notin H}\frac{P^W\p_{ov_i}(P^W\p_{ov_i})^*}{\p_{ov_i}^*Q_{v_i}\p_{ov_i}}
\right)^{-1}.
\]
\end{enumerate}
\end{theorem}

\begin{proof}
Write $\bar L=\bar L(T,\p)$ and $\bar L^{v_i}=\bar L(T^{v_i},\p^{v_i})$ for $i\in[m]$.
For the first item, let $i\in H$. For every $q\in\ker\bar L$ there holds that $q|_{T^{v_i}}\in\ker \bar L^{v_i}$ by
the characterization of the kernel in Lemma~\ref{lem:kernel}. Hence, by Lemma~\ref{lem:forbidden_directions},
\[
\langle q(v_i),\p_{ov_i}\rangle=\langle q|_{T^{v_i}}(v_i),\p_{ov_i}\rangle=0\,,
\]
due to the assumption that $i\in H$. On the other hand, by Lemma~\ref{lem:kernel} we have
\[
\langle q(o)-q(v_i),\p_{ov_i}\rangle=0\,.
\]
Therefore, $\langle q(o),\p_{ov_i}\rangle=0$. The claim follows from Lemma~\ref{lem:forbidden_directions}.

We proceed to the second item. Let $P$ be the orthogonal projection onto $\ker\bar L$.
For $z=\ii\cdot s$ with $s>0$ real, set
\[
R(z):=(\bar L-zI)^{-1}.
\]
Define $h_z(\lambda):=-z(\lambda-z)^{-1}$ for $\lambda\in\R$. By the functional calculus,
$h_z(\bar L)=-zR(z)$ is a bounded operator. For $\lambda\neq 0$ we have $h_z(\lambda)\to 0$ as $z\to 0$, while $h_z(0)=1$.
Moreover,
\[
|h_z(\lambda)|=\frac{|z|}{|\lambda-z|}\le 1\,,
\]
since $z$ is purely imaginary. Thus $h_z\to \mathbf 1_{\{0\}}$ pointwise and $(h_z)$ is uniformly bounded, so by dominated
convergence in the spectral theorem,
\[
h_z(\bar L)f\to Pf\qquad\text{for every }f\in\cH
\]
as $z\to 0$. In particular, applying this entrywise to the $d\times d$ block at $o$ gives
\[
-zR(z)_{o,o}\to P_{o,o}=Q_o\quad \mbox{as }z\to 0.
\]

Similarly, let
\[
R'(z):=\left(\sum_{i=1}^m \bar L^{v_i}-zI\right)^{-1}.
\]
Since $T$ is a tree, $\sum_{i=1}^m\bar L^{v_i}$ is a direct sum, hence so is $R'(z)$, and therefore for each $i\in[m]$,
\[
-zR'(z)_{v_iv_i}\to Q_{v_i}
\quad \mbox{as }z\to 0.
\]

We now derive the recursion. We can write
\[
\bar L=\sum_{i=1}^m \bar L^{v_i}+M,
\]
where $M$ is the rigidity Laplacian of the star at $o$. 
More explicitly, by letting $N_i=\p_{ov_i}\p_{ov_i}^*$ for $i\in [m]$, we have
\[
M_{uv}=
\begin{cases}
\sum_{j=1}^m N_j, & (u,v)=(o,o),\\
-N_i, & (u,v)=(o,v_i)\ \text{or}\ (v_i,o),\\
N_i, & (u,v)=(v_i,v_i),\\
0, & \text{otherwise}.
\end{cases}
\]
By the tree structure,
\begin{equation}\label{eq:Rprime_blocks}
R'_{o,o}(z)=-z^{-1}I,\qquad R'_{o,v_i}(z)=0,\qquad R'_{v_i,v_j}(z)=0\ \ (i\neq j).
\end{equation}

The second resolvent identity gives $R(z)-R'(z)=-R(z)MR'(z)$ (we sometimes omit the dependence on $z$ for brevity). 
Recalling \eqref{eq:Rprime_blocks}, we obtain that for every $i\in[m]$,
\begin{align*}
R_{o,v_i}
&=-R_{o,o}M_{o,v_i}R'_{v_i,v_i}-R_{o,v_i}M_{v_i,v_i}R'_{v_i,v_i}\\
&=R_{o,o}N_iR'_{v_i,v_i}-R_{o,v_i}N_iR'_{v_i,v_i}.
\end{align*}

so
\begin{equation}\label{eq:Rovi}
R_{o,v_i}=R_{o,o}N_iR'_{v_i,v_i}\,(I+N_iR'_{v_i,v_i})^{-1}.
\end{equation}
Note that in \eqref{eq:Rovi} we used that $I+N_iR'_{v_i,v_i}$ is invertible. This holds since $N_i=\p_{ov_i}\p_{ov_i}^*$, whence matrix determinant lemma yields
\[
\det(I+N_iR'_{v_i,v_i})=1+\p_{ov_i}^*R'_{v_i,v_i}\p_{ov_i}.
\]
For $z=\ii s$ with $s>0$, since $R'$ is the resolvent of a self-adjoint operator, we have
\[
\Im\bigl(\p_{ov_i}^*R'_{v_i,v_i}\p_{ov_i}\bigr)>0.
\]
In particular, $\p_{ov_i}^*R'_{v_i,v_i}\p_{ov_i}\notin\R$,
hence the determinant above is nonzero.  

Next, using the Sherman--Morrison formula for rank-one updates, \eqref{eq:Rovi} simplifies to
\begin{equation}\label{eq:Rovi_simplified}
R_{o,v_i}
=
R_{o,o}\cdot \frac{N_iR'_{v_i,v_i}}{1+\p_{ov_i}^*R'_{v_i,v_i}\p_{ov_i}}\,.
\end{equation}

Applying again $R-R'=-RMR'$ for the $(o,o)$ block, and using \eqref{eq:Rprime_blocks}, gives
\begin{align*}
R_{o,o}+z^{-1}I
&=-R_{o,o}M_{o,o}R'_{o,o}-\sum_{i=1}^m R_{o,v_i}M_{v_i,o}R'_{o,o}\\
&=z^{-1}\left(R_{o,o}\sum_{i=1}^m N_i-\sum_{i=1}^m R_{o,v_i}N_i\right)\\
&=z^{-1}R_{o,o}\sum_{i=1}^m\left(N_i-\frac{N_iR'_{v_i,v_i}N_i}{1+\p_{ov_i}^*R'_{v_i,v_i}\p_{ov_i}}\right)\\
&=z^{-1}R_{o,o}\sum_{i=1}^m\frac{N_i}{1+\p_{ov_i}^*R'_{v_i,v_i}\p_{ov_i}}\,,
\end{align*}
where in the third equality we used \eqref{eq:Rovi_simplified}, and in the fourth equality we used
$N_iR'_{v_i,v_i}N_i=(\p_{ov_i}^*R'_{v_i,v_i}\p_{ov_i})\,N_i$.
Rearranging yields the recursion
\begin{equation}\label{eq:recursion}
-zR_{o,o}(z)\left(I+\sum_{i=1}^m\frac{\p_{ov_i}\p_{ov_i}^*}{-z+\p_{ov_i}^*(-zR'_{v_i,v_i}(z))\p_{ov_i}}\right)=I.
\end{equation}

Now fix $w\in W$. Since $\p_{ov_i}^*w=0$ for $i\in H$, multiplying \eqref{eq:recursion} by $w$ gives
\[
-zR_{o,o}(z)\left(
w+\sum_{i\notin H}\frac{\p_{ov_i}\p_{ov_i}^*w}{-z+\p_{ov_i}^*(-zR'_{v_i,v_i}(z))\p_{ov_i}}
\right)=w.
\]
Letting $z\to 0$ along the imaginary axis, we use $-zR_{o,o}(z)\to Q_o$ and $-zR'_{v_i,v_i}(z)\to Q_{v_i}$.
Moreover, for $i\notin H$ we have $\p_{ov_i}^*Q_{v_i}\p_{ov_i}\neq 0$ by definition of $H$, hence the denominators converge to
$\p_{ov_i}^*Q_{v_i}\p_{ov_i}$. Therefore,
\[
Q_o\left(
w+\sum_{i\notin H}\frac{\p_{ov_i}\p_{ov_i}^*w}{\p_{ov_i}^*Q_{v_i}\p_{ov_i}}
\right)=w\,.
\]
Finally, by the first item we have $Q_oP^W=Q_o$, and clearly $P^Ww=w$ for $w\in W$. Hence, replacing $\p_{ov_i}$ by $P^W\p_{ov_i}$ inside the
rank-one terms does not change the action on $W$, and we obtain
\[
Q_o|_W=
\left(
\Id_W+\sum_{i\notin H}\frac{P^W\p_{ov_i}(P^W\p_{ov_i})^*}{\p_{ov_i}^*Q_{v_i}\p_{ov_i}}
\right)^{-1},
\]
as claimed.
\end{proof}


In the next lemma we study the event that in the random framework $(T,\nu)$, the forward flexibility matrix of the root $o$ of $T$ is trivial.

\begin{lemma}\label{lem:Qo0}
    Let $\nu$ be an admissible probability measure on $\R^d$, and let $(T,o)$ be a rooted tree such that $(T,\nu)$ is almost surely self-adjoint. 
    Then, almost surely, $Q_o=0$ if and only if $o$ has at least $d$ children $v$ for which $Q_v=0$. 
    In particular, the event $Q_o=0$ is independent of $\p(o)$.
\end{lemma}
Note that for a fixed tree \((T,o)\), the lemma shows that \(Q_o=0\) is a tail event, and therefore admits a \(0/1\) law.
However, we do not know whether the probability of this event can depend on the admissible distribution \(\nu\), or whether it depends only on the tree structure.

\begin{proof}
First, we claim that almost surely, for every edge $uv\in E(T)$ with $u$ the parent of $v$, we have
\begin{equation}\label{eq:QvpuvQv0}
Q_v\p_{uv}=0 \quad\Longleftrightarrow\quad Q_v=0.    
\end{equation}
Indeed, $\p(u)$ is independent of $Q_v$, and by our assumption on $\nu$, the vector $\p_{uv}$ almost surely does not belong to $\ker Q_v$, unless $Q_v=0$.

Denote by $v_1,\ldots,v_m$ the children of the root $o$. By \eqref{eq:QvpuvQv0}, the set $H$ from Theorem \ref{thm:recursion} is almost surely equal to
\begin{equation}\label{eq:H}
H=\{i\in[m]\mid Q_{v_i}=0\}.
\end{equation}
Theorem~\ref{thm:recursion} implies that $Q_o=0$ if and only if the subspace $W$, which is the orthogonal complement of the linear span of $\{\p_{ov_i}: i\in H\}$, is trivial.
By our assumption on $\nu$, this occurs if and only if $|H|\ge d$, as claimed.
\end{proof}


\subsection{Galton--Watson frameworks}
\label{sec:GW}
We now turn to study the local flexibility of a Galton--Watson $\mathrm{GW}(X,Y)$ branching process $(T,o)$ as defined in Section \ref{sec:preliminaries}.
Here $X,Y$ are probability distributions on $\N_{\ge 0}$ that describe the offspring distribution of the root and every other vertex respectively.
Mostly, we will be interested in the case where $Y$ is the size-biasing of $X$. Also, we restrict to frameworks of the form $(T,\nu)$, where $\nu$ is an admissible probability measure on $\R^d$. 

First, we state a sufficient condition for the self-adjointness of such frameworks.
\begin{claim}\label{clm:GWSA}
    Let $T\sim\mathrm{GW}(X,Y)$ and $\nu$ be an admissible probability measure on $\R^d$. 
    If $\E[Y]<\infty$ then $(T,\nu)$ is almost surely self-adjoint.
\end{claim}
Note that if $Y$ is the size-biasing of $X$ then $\E[Y]<\infty$ if and only if $\E[X^2]<\infty$.
\begin{proof}
There exists $C\in\N$ such that
\begin{equation}\label{eq:chooseC}
\sum_{m>C} m\,Y(m)<1.
\end{equation}
We claim that, almost surely, for every vertex $v$ there exists a finite subtree $F^v$ of $T$, rooted at $v$, such that every leaf of $F^v$
has at most $C$ children in $T$. Clearly, it suffices to prove this for every non-root vertex $v$.

Reveal the tree $T$ from $v$ outward, but whenever an explored vertex has at most $C$ children, declare it a leaf and do not explore its children.
Let $F^v$ be the subtree revealed by this process. Then $F^v$ is a Galton--Watson tree rooted at $v$ with offspring distribution $Y'$ given by
\[
Y'(m)=
\begin{cases}
\P(Y\le C), & m=0,\\
0, & 0<m\le C,\\
Y(m), & m>C,
\end{cases}
\]
with expected number of offspring equal to
\[
\E[Y']=\sum_{m>C} m\,Y(m)<1\,.
\]
Hence, $F^v$ is almost surely finite. Moreover, by construction, every leaf of $F^v$ has at most $C$ children in $T$.

Let $S_r\subseteq V(T)$ consist of $o$ and the union of $F^v$ over all vertices $v$ at distance at most $r$ from $o$.
Then $S_r$ induces a connected finite subtree, contains the root $o$, and it is $C$-bounded. In addition, $\bigcup_{r\ge 0} S_r=V(T)$.
Since $\nu$ has a bounded support, the claim follows from Lemma~\ref{lem:tree_self_adjoint}.
\end{proof}


The main result of this section solves the recursion in Theorem~\ref{thm:recursion} for Galton--Watson frameworks in which $Y$ is the size-biasing of $X$,
and the embeddings are sampled independently from an admissible distribution $\nu$.

\begin{theorem}\label{thm:Exp_loc_flex}
    Suppose that $X$ is a probability distribution on $\N_{\ge 0}$ with a bounded second moment, $Y$ the size-biasing of $X$, $T\sim \mathrm{GW}(X,Y)$ with root $o$, and let $\nu$ be a probability measure on $\R^d$ as above. Then, there exists a solution $p \in [0,1]$ to $p=\P(Y_p\ge d)$ such that
    \begin{equation}\label{eq:Xformula}
    \E[\phi_{T,\nu}(o)] = \E[\mathbf 1_{X_p<d}\cdot(d-X_{p})] - \frac12(1-p)^2\E[X]\,.
    \end{equation}
    
\end{theorem}
\begin{proof}
Let $(T',u)$ denote a $\mathrm{GW}(Y,Y)$-distributed tree rooted at $u$. Then, $(T',\nu)$ is almost surely self-adjoint by Claim \ref{clm:GWSA}. 
Let $p_Y=\P(Q_u=0)$, where $Q_u$ denotes the forward flexibility matrix of $u$ in $T'$. 
We claim that the number of children $v$ of $u$ in $T'$ for which $Q_v=0$ is $Y_{p_Y}$-distributed. 
Indeed, the total number of children $v$ of $u$ is $Y$-distributed, and the events $Q_v=0$ are independent, each with probability $p_Y$, since the subtree rooted at $v$ is also $\mathrm{GW}(Y,Y)$-distributed.
By Lemma \ref{lem:Qo0} we find that 
\begin{equation}\label{eq:px-py-bin-recursion}
p_Y=\P\!\left(Y_{p_Y}\ge d\right)\,.
\end{equation}

If \(p_Y=1\), then \eqref{eq:px-py-bin-recursion} gives
\(\mathbb P(Y\ge d)=1\), and hence
    $\mathbb P(1\le X\le d)=0$.
Substituting \(p=1\) into the right-hand side of
\eqref{eq:Xformula} therefore gives \(dX(0)\).
On the other hand, the root \(o\) of \(T\) is isolated with probability
\(X(0)\). In this case its local flexibility is \(d\). Otherwise, \(o\)
has at least \(d+1\) children, and all their forward flexibility matrices
vanish. Hence, the local flexibility at the root is \(0\) by
Lemma~\ref{lem:Qo0}. Thus the expected local flexibility is \(dX(0)\),
so \eqref{eq:Xformula} holds also in this case.

So we assume $p_Y<1$. The following distribution plays a key role in the proof: sample a rooted $\mathrm{GW}(Y,Y)$-tree $(T',u)$ and an embedding $\p:V(T')\to\R^d$ by drawing i.i.d.\ samples from $\nu$. Let $\cD^Y$ denote the law of the pair $(\p(u),Q_u)$ under this procedure, conditioned on the event $Q_u\neq 0$.
By Lemma~\ref{lem:Qo0}, the marginal law of $\p(u)$ under $\cD^Y$ is $\nu$, since $\p(u)$ is independent of the event $Q_u=0$. 
However, the matrix $Q_u$ itself may depend on $\p(u)$.

\medskip


We now turn our attention to $(T,o)$. Note that $(T,\nu)$ is almost surely self-adjoint by Claim \ref{clm:GWSA}. 
By \eqref{eq:QvpuvQv0}, the subset $H$ from Theorem \ref{thm:recursion} is almost surely equal to the subset of the children $v$ of the root $o$ for which $Q_v=0$. We claim that $|H|$ is $X_{p_Y}$-distributed. Indeed, the total number of children $v$ of $o$ is $X$-distributed, and, as in $T'$, the events $Q_v=0$ are independent of probability $p_Y$. 

Let $(\p(u_i),Q_{u_i})_{i\ge 1}$ be a sequence of i.i.d.\ samples from $\cD^Y$, and $\p(o),\p(v_i),~i=1,\ldots,d-1$ independently drawn from $\nu$. Via Theorem \ref{thm:recursion}, we can sample a matrix $Q$ with the same distribution as the forward flexibility matrix $Q_o$ of $o$ in $(T,\nu)$
 using these random vectors and matrices: 
Indeed, start by sampling \((h,s)\sim (X_{p_Y},X_{1-p_Y})\).
Viewing \(h\) as the size of \(|H|\), we set \(Q=0\) if \(h\ge d\).
Otherwise, view
\begin{enumerate}
    \item[(i)] \(v_1,\ldots,v_h\) as the children of the root \(o\)
    for which \(Q_{v_i}=0\) (note that \(\p(v_i)\) is
    \(\nu\)-distributed even under the conditioning \(Q_{v_i}=0\)),
    \item[(ii)] \(u_1,\ldots,u_s\) as the remaining children of the root,
    and
    \item[(iii)] \(Q_{u_i}\) as the forward flexibility matrix of \(u_i\)
    in \(T\) (matching the definition of \(\cD^Y\)).
\end{enumerate}



In this realization, $W$ is the orthogonal complement of
\[
\mathrm{span}\{\p_{ov_i}: 1\le i\le h\},
\]
and has dimension $d-h$. Set 
\[
A_i=\frac{P^W\p_{ou_i}(P^W\p_{ou_i})^*}{\p_{ou_i}^*Q_{u_i}\p_{ou_i}}\,,\qquad i\ge 1\,,
\]
which is a $(d-h)\times(d-h)$ matrix acting on $W$. Theorem~\ref{thm:recursion} implies that the random $d\times d$ matrix $Q$ defined by $Q|_{W^\perp}=0$ and
\[
Q|_W=\Bigl(\Id_W+\sum_{i=1}^{s}A_i\Bigr)^{-1}
\]
has the distribution of the forward flexibility matrix $Q_o$ of the root $o$ in $T$.
Therefore,

\begin{align}
\notag 
\E[\tr\,Q_o] 
&=~ \E\left[ \mathbf 1_{h<d}\cdot \tr\left(\left(\Id_W + \sum_{i=1}^{s}A_i\right)^{-1}\right)\right] \\
\notag &=~ \E\left[\mathbf 1_{h<d}\cdot \tr\left(
\Id_W-\sum_{i=1}^{s}A_i \left(\Id_W + \sum_{i'\in[s]\setminus\{i\}}A_{i'} +A_i\right)^{-1}\right)\right] \\
\label{eq:first_stop} &=~ \E[\mathbf 1_{h<d}\cdot (d-h)] - \E_{s\sim X_{1-p_Y}}\left[ s\cdot f(s) \right]\,,
%\sum_{i=1}^{s}\tr\left[A_i \left(\Id_W +  \sum_{j\ne s}A_j+A_i\right)^{-1}\right]\right]   \\
%\label{eq:first_stop} &=~ d-h - \E[s\cdot f_h(s)]\,,
\end{align}
where $f:\N\to \R$ is defined by
\[
f(s) = 
\E\left[\mathbf{1}_{h<d}\cdot 
\tr\left(A_s \left(\Id_W + \sum_{i=1}^{s-1}A_i+A_s\right)^{-1}\right)
~~\middle |~~s\right]\,.
\]
In the align, the transition to the second line is by the equality $(I+S)^{-1}=I-S(I+S)^{-1}$ which holds if $I+S$ is invertible, and the next transition is by linearity, $\dim W=d-h$, and the fact that the $A_i$'s have the same distribution. 

By \eqref{eq:sizeBiasF}, the equality $\E[X_{1-p_Y}]=(1-p_Y)\E[X]$, and the first item in Claim \ref{clm:sizeBias_new_claim}, we reformulate \eqref{eq:first_stop} to
\begin{equation}\label{eq:ft}
\E[\tr\,Q_o]=\E[\mathbf 1_{X_{p_Y}<d}\cdot (d-X_{p_Y})]- 
(1-p_Y)\E[X]\cdot\E_{t\sim Y_{1-p_Y}}\left[f(t+1) \right]\,.    
\end{equation}

We turn to study the expectation of $f(t+1)$, where $t\sim Y_{1-p_Y}$. 
By definition, when evaluating $f(t+1)$ we draw
$h\sim X_{p_Y}\mid \{X_{1-p_Y}=t+1\}$. Therefore, by the third item of
Claim~\ref{clm:sizeBias_new_claim}, the joint distribution of $(h,t)$ is
$(Y_{p_Y},Y_{1-p_Y})$. Once $h$ is sampled, both $W$ and the matrices $A_i$ are defined via the randomness $(\p(u_i),Q_{u_i})_{i\ge 1}$ together with the independent samples $\p(o)$ and $\p(v_i)_{i<d}$. We denote by $\E'$ the expectation taken over $(h,t)\sim(Y_{p_Y},Y_{1-p_Y})$ and all this auxiliary randomness, and obtain
\begin{equation}\label{eq:Exp_of_ft}
    \E_{t\sim Y_{1-p_Y}}\!\left[f(t+1) \right]
    = \E'\!\left[ 
        \mathbf{1}_{h<d}\cdot 
\tr\!\left(A_{t+1} \left(\Id_W + \sum_{i=1}^{t}A_i+A_{t+1}\right)^{-1}\right)
    \right].
\end{equation}

Condition on the value of $h<d$. In such a case, we let $B_t$ be the matrix of order $d-h$ acting on $W$ defined by 
$$B_t=\Id_W + \sum_{j=1}^{t}A_j\,.$$ Also, denote $u=u_{t+1}$, $Q_u=Q_{u_{t+1}}$ for brevity. By the Sherman--Morrison formula and the definition of $A_{t+1}$, we find that
\[
A_{t+1}(B_t+A_{t+1})^{-1} = \frac{A_{t+1}B_t^{-1}}{1+\frac{(P^W\p_{ou})^*B_t^{-1}(P^W\p_{ou})}{\p_{ou}^*Q_{u}\p_{ou}}}\,.
\]
Therefore, the cyclic property of the trace and a simple manipulation give
\begin{align}
\tr\!\left[A_{t+1}(B_t+A_{t+1})^{-1}\right]
&= \frac{(P^W\p_{ou})^*B_t^{-1}(P^W\p_{ou})}{\p_{ou}^*Q_{u}\p_{ou}+(P^W\p_{ou})^*B_t^{-1}(P^W\p_{ou})} \nonumber\\[0.6em]
&= \frac{\p_{ou}^*\tilde Q\,\p_{ou}}{\p_{ou}^*Q_{u}\p_{ou}+\p_{ou}^*\tilde Q\,\p_{ou}}\,, \label{eq:trace_A_over_BplusA}
\end{align}
where $\tilde Q$ denotes a $d\times d$ matrix defined by $\tilde Q|_{W^\perp}=0$ and $\tilde Q|_W = B_t^{-1}$.

By combining \eqref{eq:Exp_of_ft} and \eqref{eq:trace_A_over_BplusA} we find
\begin{equation}\label{eq:almost_there}
\E_{t\sim Y_{1-p_Y}}[f(t+1)]=\E'(\mathbf 1_{h<d})\E'\left[ \frac{\p_{ou}^*\tilde Q\,\p_{ou}}{\p_{ou}^*Q_{u}\p_{ou}+\p_{ou}^*\tilde Q\,\p_{ou}} ~~\middle |~~ h<d \right]\,.  
\end{equation}

Crucially, observe that, by combining Theorem \ref{thm:recursion} and the definition of $B_t$, we can view $\tilde Q$ as a forward flexibility matrix. More precisely, since $(h,t)\sim (Y_{p_Y},Y_{1-p_Y})$, we derive that, conditioned on $h<d$, the pair $(\p(o),\tilde Q)$ is $\cD^Y$-distributed.
In addition, it is independent of the $\cD^Y$-distributed pair $(\p(u),Q_u)$ (which by definition denotes the pair $(\p(u_{t+1}),Q_{u_{t+1}})$).
Therefore, conditioned on $h<d$, the pair of positive random variables
$\p_{ou}^*Q_{u}\p_{ou}$ and 
$\p_{ou}^*\tilde Q\,\p_{ou}$
are exchangeable: indeed, the pairs \((\p(u),Q_u)\) and \((\p(o),\tilde Q)\) are independent and $\cD^Y$-distributed, and exchanging them only replaces \(\p_{ou}\) by \(-\p_{ou}\). Hence,
\[
\E'\left[ \frac{\p_{ou}^*\tilde Q\,\p_{ou}}{\p_{ou}^*Q_{u}\p_{ou}+\p_{ou}^*\tilde Q\,\p_{ou}} ~~\middle |~~ h<d \right]=\frac{1}{2}\,.
\]
Combined with $\E'[\mathbf 1_{h<d}]=1-p_Y$ which follows from \eqref{eq:px-py-bin-recursion}, we find that \eqref{eq:almost_there} reformulates to
\[
\E_{t\sim Y_{1-p_Y}}[f(t+1)]=\frac{1-p_Y}{2}\,.
\]
Substituting this equality into \eqref{eq:ft} yields
\begin{align*}
    \E[\tr\,Q_o]
    &=~
\E[\mathbf 1_{X_{p_Y}<d}\cdot(d-X_{p_Y})] - \frac12(1-p_Y)^2\E[X]\,,
\end{align*}
and together with~\eqref{eq:px-py-bin-recursion}, the theorem follows.

\end{proof}

\begin{thebibliography}{99}
\bibitem[]{AbertThomViragSpectralMeasure} Ab{\'e}rt, Mikl{\'o}s; Thom, Andreas; Vir{\'a}g, B{\'a}lint. Benjamini--Schramm convergence and pointwise convergence of the spectral measure. (2013)
\bibitem[]{AldousSteeleObjectiveMethod} Aldous, David; Steele, J. Michael. The Objective Method: Probabilistic Combinatorial Optimization and Local Weak Convergence. Probability on Discrete Structures 110 (2004) 1--72
\bibitem[]{BenjaminiSchramm2001} Benjamini, Itai; Schramm, Oded. Recurrence of Distributional Limits of Finite Planar Graphs. Electronic Journal of Probability 6 (2001) 1--13
\bibitem[]{BordenaveLelarge2010} Bordenave, Charles; Lelarge, Marc. Resolvent of Large Random Graphs. Random Structures \& Algorithms 37 (2010) 332--352
\bibitem[]{BordenaveLelargeSalezDiluted} Bordenave, Charles; Lelarge, Marc; Salez, Justin. The Rank of Diluted Random Graphs. The Annals of Probability 39 (2011) 1097--1121
\bibitem[]{KitsonPowerRigidityInfiniteGraphs} Kitson, Derek; Power, Stephen C.. The rigidity of infinite graphs. Discrete \& Computational Geometry 60 (2018) 531--557
\bibitem[]{OwenPowerInfiniteFrameworksOperatorTheory} Owen, John C.; Power, Stephen C.. Infinite bar-joint frameworks, crystals and operator theory. New York Journal of Mathematics 17 (2011) 445--490
\bibitem[]{ReedSimonI} Reed, Michael; Simon, Barry. Methods of Modern Mathematical Physics. {I}: Functional Analysis. (1980)
\end{thebibliography}
\end{document}
